\section{State Machine}
\label{sec:state-machine}

\subsection{Overview}
The \texttt{StateMachine} class in \texttt{MLGWorks.Utils.Patterns} provides a lightweight and extensible framework for implementing finite state machines. It allows you to define states via the \texttt{IState} interface and manage transitions and ticking logic.

\subsection{Interfaces}

\begin{verbatim}
public interface IState
{
    void OnEnter();
    void OnExit();
    void Tick();
}
\end{verbatim}

Any class implementing \texttt{IState} can be used as a state in the state machine. The interface defines three methods:
\begin{itemize}
    \item \texttt{OnEnter()}: Called when the state becomes active.
    \item \texttt{OnExit()}: Called when the state is deactivated.
    \item \texttt{Tick()}: Called once per frame while the state is active.
\end{itemize}

\subsection{Usage}

\begin{verbatim}
public class IdleState : IState
{
    public void OnEnter() => Debug.Log("Entering Idle");
    public void OnExit() => Debug.Log("Exiting Idle");
    public void Tick() => Debug.Log("Idling...");
}

public class PatrolState : IState
{
    public void OnEnter() => Debug.Log("Start Patrol");
    public void OnExit() => Debug.Log("Stop Patrol");
    public void Tick() => Debug.Log("Patrolling...");
}
\end{verbatim}

\textbf{Creating and using the state machine:}

\begin{verbatim}
StateMachine fsm = new StateMachine();

IState idle = new IdleState();
IState patrol = new PatrolState();

fsm.ChangeState(idle);  // Enters idle
fsm.Tick();             // Calls idle.Tick()

fsm.ChangeState(patrol); // Exits idle, enters patrol
fsm.Tick();              // Calls patrol.Tick()
\end{verbatim}

\subsection{Features}
\begin{itemize}
    \item Simple and intuitive API for switching states.
    \item Clean separation of logic per state.
    \item Ensures proper lifecycle management with \texttt{OnEnter} and \texttt{OnExit}.
    \item Decouples state logic from the state machine controller.
\end{itemize}

\subsection{Best Practices}
\begin{itemize}
    \item Use this system for AI, gameplay logic, input handling, and other behavior switching tasks.
    \item Consider combining with a visual editor or scriptable objects for more advanced workflows.
    \item Extend with transition conditions, logging, or debugging utilities as needed.
\end{itemize}
